\newpage
\section{Soal 16: MLP untuk Prediksi Horsepower (Auto MPG)}

Dataset D8 Auto MPG berisi 398 mobil keluaran 1970 sampai 1982. Model dibangun dengan \texttt{PyTorch}.

\subsection{Business Understanding}
Data spesifikasi kendaraan sering bolong, termasuk horsepower yang dipakai untuk pajak, asuransi, dan perbandingan performa. Kalau horsepower bisa ditebak akurat dari spesifikasi lain, data yang bolong bisa dilengkapi. Seperti montir berpengalaman yang bisa menebak tenaga mesin dari ukuran mesin dan bobot mobil.

\begin{tabularx}{\textwidth}{@{}LL@{}}
\toprule
\textbf{Kriteria kesuksesan} & \textbf{Ambang} \\
\midrule
$R^2$ test & $\geq 0.85$ \\
MAE test & $\leq 10$ hp \\
Tidak overfit berat & RMSE test $\leq 1.3 \times$ RMSE train \\
Lebih baik dari atau setara regresi linear & RMSE test MLP $\leq$ linear \\
\bottomrule
\end{tabularx}

\subsection{Data Understanding}
Dari 398 mobil tidak ada duplikat. Missing value hanya ada di horsepower, sebanyak 6 baris: Ford Pinto, Ford Maverick, Renault LeCar, Ford Mustang Cobra, Renault 18i, dan AMC Concord. Horsepower berkisar 46 sampai 230 hp (rata-rata 104, std 38.5) dan miring ke kanan. Ke-10 outlier IQR semuanya mobil V8 bertenaga besar, bukan salah input.

Horsepower berhubungan kuat dengan fitur lain: korelasinya 0.90 dengan displacement, 0.87 dengan weight, 0.84 dengan cylinders, $-0.78$ dengan mpg, dan $-0.69$ dengan acceleration. Scatter horsepower terhadap mpg melengkung, jadi MLP yang bisa menangkap pola non-linear layak dicoba.

\fig{soal_16_1.png}{0.95}{Distribusi horsepower dan hubungannya dengan fitur lain (warna = origin).}

\jawab{Jawaban 16a: Preprocessing dan urutan agar tidak bocor}
\begin{tabularx}{\textwidth}{@{}lp{2.6cm}L@{}}
\toprule
\textbf{Langkah} & \textbf{Kapan / dari mana} & \textbf{Alasan} \\
\midrule
1. Hapus 6 baris horsepower kosong & sebelum split (sesuai soal) & horsepower adalah target, dan mengisinya dengan mean sama dengan mengarang kunci jawaban. Menghapus tidak menghitung statistik apa pun, jadi aman dilakukan sebelum split. \\
2. Buang \texttt{car\_name} & sebelum split & sesuai soal; teks unik per mobil \\
3. Split 70/15/15 (seed 47) & - & 274 train, 59 validation, 59 test \\
4. One-hot \texttt{origin} (USA/Eropa/Jepang) & kategori tetap & angka 1/2/3 hanya kode; kalau dibiarkan numerik, model mengira Jepang = 3$\times$ USA \\
5. StandardScaler 6 fitur numerik & fit di train saja & MLP sensitif skala (berat sekitar 3000, silinder sekitar 5); statistik test tidak boleh dipakai \\
6. Standardisasi target & mean/std train & gradien lebih stabil; prediksi dikembalikan ke satuan hp sebelum metrik dihitung \\
\bottomrule
\end{tabularx}

\emph{Data leakage} mirip siswa yang sudah melihat soal ujian saat belajar: nilainya bagus, tetapi tidak menggambarkan pemahaman. Outlier tidak saya buang karena mobil V8 bertenaga besar memang ada, dan model perlu belajar dari mereka.

\jawab{Jawaban 16b: Rancangan MLP}
Arsitekturnya \texttt{input(9) $\to$ Linear(64) $\to$ ReLU $\to$ Linear(32) $\to$ ReLU $\to$ Linear(1)} dengan 2.753 parameter.

\begin{tabularx}{\textwidth}{@{}llL@{}}
\toprule
\textbf{Komponen} & \textbf{Pilihan} & \textbf{Alasan} \\
\midrule
Hidden layer & 64 dan 32 neuron & memenuhi syarat minimal 2 hidden layer; bentuk corong dengan ukuran wajar untuk 274 data \\
Aktivasi & ReLU & cepat, gradien tidak menghilang, bisa menangkap pola melengkung \\
Output & 1 neuron linear & regresi, jadi output berupa angka bebas \\
Loss & MSE & standar regresi dan sejalan dengan RMSE \\
Optimizer & Adam, lr 0.001, batch 32 & menyesuaikan langkah tiap bobot, stabil untuk data kecil \\
Early stopping & patience 100, maks 3000 epoch & ambil bobot dengan RMSE validation terbaik \\
\bottomrule
\end{tabularx}

\jawab{Jawaban 16c: Hasil, \emph{learning curve}, dan prediksi vs aktual}
\begin{adjustbox}{max width=\textwidth}\begin{tabular}{@{}lccc@{}}
\toprule
\textbf{Data} & \textbf{MAE (hp)} & \textbf{RMSE (hp)} & \textbf{$R^2$} \\
\midrule
MLP train & 7.48 & 10.63 & 0.927 \\
MLP validation & 6.30 & 9.05 & 0.948 \\
MLP test & 5.40 & 7.10 & 0.951 \\
Regresi linear test & 6.51 & 7.82 & 0.941 \\
\bottomrule
\end{tabular}\end{adjustbox}

\fig{soal_16_2.png}{1.0}{Kiri: \emph{learning curve} dengan bobot terbaik di epoch 28. Kanan: prediksi vs aktual di test set.}

Di test set, MLP rata-rata meleset 5.4 hp dan menjelaskan 95.1\% variasi horsepower. MAE-nya sekitar 17\% lebih kecil dari regresi linear. Setelah kira-kira epoch 25, RMSE validation mendatar di sekitar 9 hp sementara RMSE train terus turun, tanda model mulai menghafal. Early stopping mengambil bobot epoch 28, sebelum hafalan itu terjadi. Penyebaran prediksi terbesar ada pada mobil V8 140 sampai 200 hp.

\jawab{Jawaban 16d: Overfit atau underfit, dan eksperimen satu perubahan}
Model ini tidak overfit. RMSE train (10.6) malah lebih besar dari validation (9.0) dan test (7.1). Porsi mobil di atas 202 hp hampir sama di train (2.9\%) dan validation (3.4\%), jadi penyebabnya bukan data latih yang lebih sulit. Ada tiga hal yang menjelaskannya. Saat bobot terbaik diambil di epoch 28, model belum selesai mempelajari data latih. RMSE validation sedikit optimis karena epoch terbaik justru dipilih berdasarkan validation. Test set kebetulan lebih mudah: tidak ada satu pun mobil di atas 202 hp, dan std-nya 32.5 dibanding 39.5 di train.

Untuk eksperimen, saya mengubah satu elemen saja: menambah Dropout 0.2 setelah tiap hidden layer. Dropout mematikan 20\% neuron secara acak selama latihan, mirip belajar kelompok yang anggotanya bergantian absen sehingga semua terpaksa paham materinya. Kedua varian diulang dengan 5 seed:

\begin{adjustbox}{max width=\textwidth}\begin{tabular}{@{}lcccccc@{}}
\toprule
\textbf{Varian} & \textbf{RMSE train} & \textbf{RMSE val} & \textbf{RMSE test} & \textbf{MAE test} & \textbf{$R^2$ test} & \textbf{Epoch terbaik} \\
\midrule
Awal & $9.71 \pm 0.74$ & $9.06 \pm 0.17$ & $7.35 \pm 0.37$ & $5.71 \pm 0.34$ & 0.948 & 47 \\
Dropout 0.2 & $8.99 \pm 0.79$ & $8.86 \pm 0.38$ & $6.97 \pm 0.30$ & $5.48 \pm 0.25$ & 0.953 & 128 \\
\bottomrule
\end{tabular}\end{adjustbox}

\fig{soal_16_3.png}{0.95}{Learning curve model awal vs Dropout 0.2 (seed 47).}

Dropout menurunkan RMSE validation, yang menjadi dasar pemilihan, dan RMSE test ikut turun. Latihan menjadi lebih berisik sehingga model tidak cepat menghafal, dan bobot terbaik baru tercapai di sekitar epoch 128, jauh setelah epoch 47 pada model awal. Perbaikannya kecil, 0.2 sampai 0.4 hp, kira-kira sebesar satu simpangan baku antar-seed. Karena itu eksperimen ini saya ulang 5 seed; dengan satu seed saja kesimpulannya bisa terbalik.

\subsection{Interpretasi Hasil}
Pada 20 mobil test, kebanyakan prediksi meleset kurang dari 7 hp. Ford Galaxie 500, misalnya, punya horsepower 198 dan diprediksi 197.9. Error terbesar ada pada Ford LTD (+15 hp). Test set tidak memuat mobil di atas 202 hp, jadi kemampuan model untuk mobil bertenaga sangat besar belum benar-benar teruji. Untuk enam mobil tanpa label, semua prediksi masuk rentang horsepower mobil latih dengan jumlah silinder yang sama; Renault LeCar diprediksi 57 hp dan Ford Mustang Cobra 101 hp.

\begin{adjustbox}{max width=\textwidth}\begin{tabular}{@{}llll@{}}
\toprule
\textbf{Kriteria} & \textbf{Target} & \textbf{Hasil} & \textbf{Status} \\
\midrule
$R^2$ test & $\geq 0.85$ & 0.951 & \ok{tercapai} \\
MAE test & $\leq 10$ hp & 5.40 hp & \ok{tercapai} \\
Tidak overfit berat & $\leq 1.3\times$ train & 7.10 $\leq$ 13.8 & \ok{tercapai} \\
$\leq$ regresi linear & RMSE & 7.10 $<$ 7.82 & \ok{tercapai} \\
\bottomrule
\end{tabular}\end{adjustbox}

Dengan pipeline tanpa kebocoran, MLP 9-64-32-1 mencapai MAE 5.4 hp dan $R^2$ 0.951 tanpa overfit, berkat early stopping. Dropout 0.2 memberi perbaikan kecil tetapi konsisten. Datanya kecil, jadi \emph{k-fold cross-validation} akan memberi estimasi performa yang lebih stabil.
